\begin{table*}[t!]
\centering
\begin{threeparttable}
\caption{Cross-Engine Empirical Performance Comparison Between Serverless Cloud OLAP (Google BigQuery) and In-Process Vectorized OLAP (DuckDB) Across Five SQL Complexity Tiers (TPC-H Scale Factor = 1 GB)}
\label{tab:cross_engine_main}
\setlength{\tabcolsep}{5.2pt}
\renewcommand{\arraystretch}{1.18}
\begin{tabular}{@{}lccccccccc@{}}
\toprule
\multirow{2}{*}{\textbf{Complexity Tier}} & \multicolumn{2}{c}{\textbf{AST Structure}} & \multicolumn{4}{c}{\textbf{Google BigQuery Serverless (Cloud)}} & \multicolumn{2}{c}{\textbf{DuckDB Vectorized (Local)}} & \multirow{2}{*}{\textbf{Ratio}} \\
\cmidrule(lr){2-3} \cmidrule(lr){4-7} \cmidrule(lr){8-9}
 & \textbf{Depth} & \textbf{Joins} & \textbf{Unopt (ms)} & \textbf{Clustered (ms)} & \textbf{Slot (ms)} & \textbf{Scanned (MB)} & \textbf{Parquet (ms)} & \textbf{Native (ms)} & \textbf{(BQ/Duck)} \\
\midrule
\textbf{L1: Scan \& Filter} & $9.0$ & $0.0$ & $1368 \pm 2111$ & $433 \pm 658$ & $1981$ & $332.5 \rightarrow 166.8$ & $81.6 \pm 110.7$ & $\mathbf{50.8 \pm 90.5}$ & $8.5\times$ \\
\textbf{L2: Multi-Way Joins} & $9.2$ & $5.2$ & $3667 \pm 3660$ & $667 \pm 824$ & $5194$ & $473.9 \rightarrow 333.3$ & $403.9 \pm 521.6$ & $\mathbf{264.5 \pm 378.7}$ & $2.5\times$ \\
\textbf{L3: Nested CTEs} & $10.8$ & $4.6$ & $1700 \pm 1835$ & $493 \pm 265$ & $4637$ & $417.3 \rightarrow 259.4$ & $232.3 \pm 264.1$ & $\mathbf{121.9 \pm 140.9}$ & $4.0\times$ \\
\textbf{L4: Window Func.} & $10.8$ & $1.2$ & $1012 \pm 1146$ & $373 \pm 455$ & $1908$ & $241.5 \rightarrow 123.2$ & $111.5 \pm 80.6$ & $\mathbf{78.2 \pm 70.2}$ & $4.8\times$ \\
\textbf{L5: Array / UNNEST} & $10.6$ & $3.0$ & $3206 \pm 3280$ & $1761 \pm 1576$ & $8737$ & $273.4 \rightarrow 180.6$ & $442.7 \pm 683.6$ & $\mathbf{403.3 \pm 640.0}$ & $4.4\times$ \\
\bottomrule
\end{tabular}
\begin{tablenotes}\footnotesize
\item \textit{Note:} All measurements were executed with query caching disabled (\texttt{use\_query\_cache=False}) after warm-up runs. Scanned (MB) indicates transitions from unoptimized tables to clustered native tables. Acknowledgment: This research project was financially supported by Mahasarakham University.
\end{tablenotes}
\end{threeparttable}
\end{table*}
